\chapter{Reference}
\label{ch:reference}

\section{Conclusions, stated compactly}

\begin{description}[leftmargin=0cm,style=nextline,itemsep=6pt]

\item[What reproduces exactly.] The checkpoint is the model described
  (253{,}556{,}736 parameters). It returns the published property vector for the
  published worked example --- for both printed variants of it. The 1{,}033-pair
  dataset reconstructs to the exact published count. The normalisation constants
  were recovered independently and later confirmed against the real Table~S5,
  all eight. Given the paper's own designed sequences, the forward task returns
  the published \Rsq{} to four decimals on four of five sets, matching the whole
  predicted vector slot for slot. Descriptors match on 15/15 molecular weights
  and isoelectric points.

\item[What reproduces partially.] Running the design procedure end to end, my
  maximum falls below the published value on all eight targets, mean gap 0.27.
  Because the scorer reproduces exactly, this is located in the inverse task's
  \emph{search}, not in the model. Why the authors' search found longer and
  better candidates, I cannot say.

\item[What the numbers measure.] The headline \Rsq{} is a maximum over a
  sampling budget that is not stated, of a within-vector self-consistency
  measure, from a model that returns the exact training label for 71\% of its
  fine-tuning pairs and emits verbatim training sequences 62\% of the time.
  Each qualifier is individually reasonable; together they mean the number
  measures something narrower than it appears to.

\item[What would remove the ambiguity.] Two changes, neither requiring new
  experiments: state $N$ alongside any best-of-$N$ result, and hold out a fold.
  On 1{,}033 pairs a held-out fold is cheap, and the baselines suggest it would
  be informative --- the gap between in-sample and cross-validated performance on
  this data is the difference between 0.80 and $-0.16$.

\item[The wider point.] Generative design pipelines that chain an inverse task
  into a forward task and report the agreement are measuring something real.
  But what they measure depends on the search budget and on how much of the
  training set the model has retained. Both are cheap to report and usually are
  not.

\end{description}

\section{The fifteen assumptions}

Every choice the paper does not specify. Full reasoning in
\code{ASSUMPTIONS.md}.

\begin{longtable}{@{}p{0.7cm}p{6.0cm}p{7.4cm}@{}}
\toprule
\textbf{ID} & \textbf{the choice} & \textbf{basis} \\
\midrule
\endhead
A1 & Dataset reconstruction rule & Join on \code{idv\_id}; lands on 1{,}033 exactly \\
A2 & Normalisation constants & Recovered; verified 0.00047, later confirmed vs Table~S5 \\
A3 & Validity filter on generations & Twenty standard residues only; rejection rate reported \\
A4 & Definition of novelty & Paper's exact-match test, plus $k$-mer measures; calibrated in §\ref{sec:exactmatch} \\
A5 & Greedy forward decoding & 58/60 identical to sampled; batch-invariant \\
A6 & Inference precision & float16 on GPU; determinism verified at that precision \\
A7 & Terminal-domain boundaries & 130 / 100 residues, nominal literature values \\
A8 & Motif set & Superseded: Table~S4 obtained, all 14 entries verified \\
A9 & Sampling budget $N$ & Treated as a variable; the whole curve reported \\
A10 & Batched inference & Equivalence to unbatched tested \\
A11 & Checkpoint revision & Resolved commit recorded and pinnable \\
A12 & Duplicate sequence labels & Kept; the published count requires it \\
A13 & Composition accessor & Counted directly; avoids a Biopython API change \\
A14 & Secondary-structure convention & Both legacy and corrected reported (see D1) \\
A15 & Amino-acid grouping & Two schemes reported; neither chosen to fit \\
\bottomrule
\end{longtable}

\section{The fifteen discrepancies and open points}

\begin{longtable}{@{}p{0.7cm}p{6.0cm}p{7.4cm}@{}}
\toprule
\textbf{ID} & \textbf{point} & \textbf{status} \\
\midrule
\endhead
D1 & Figure 5c helix/sheet labels exchanged vs current Biopython & confirmed; upstream correction postdates publication \\
D2 & Set S3 strength 0.200 where the stated median is 0.310 & confirmed, internal to the paper \\
D3 & Sampling budget $N$ never stated & clarification; quantified in §\ref{sec:bestofn} \\
D4 & \Rsq{} is within-vector, not across a test set & clarification \\
D5 & No held-out set & clarification; stated by the paper itself \\
D6 & 62.3\% of generations are verbatim training sequences & confirmed \\
D7 & Supporting Information initially unobtainable & resolved via arXiv preprint \\
D8 & Paper prints two variants of one worked example & confirmed; my earlier self-blame retracted \\
D9 & \code{r2\_score} argument order differs between two helpers & clarification; no published number affected \\
D10 & Shortfall on all eight targets & confirmed; located in the search (§\ref{sec:exactmatch}) \\
D11 & Forward task returns 71\% of training labels verbatim & confirmed \\
D12 & A reviewer comment described a different paper's figure & clarification \\
D13 & \textbf{Withdrawn.} Alleged motif sign discrepancy & \textbf{retracted}; all 14 entries agree \\
D14 & Recovered constants match published Table~S5 exactly & confirmed \\
D15 & Table~S1 extraction & failed, then rewritten and verified \\
\bottomrule
\end{longtable}

\section{Where every number comes from}

\begin{longtable}{@{}p{6.6cm}p{7.5cm}@{}}
\toprule
\textbf{script} & \textbf{produces} \\
\midrule
\endhead
\code{00\_check\_env.py} & architecture check, pinned revision \\
\code{01\_smoke\_test.py} & both tasks end to end, cost projection \\
\code{01b\_forward\_decoding\_diagnostics.py} & the evidence for assumption A5 \\
\code{02\_reproduce\_table1.py} & the eight candidate pools (the long run) \\
\code{03\_build\_dataset.py} & the 1{,}033 pairs; normalisation check \\
\code{04\_forward\_eval\_dataset.py} & in-sample forward accuracy \\
\code{04b\_memorisation\_audit.py} & the 71.3\% figure \\
\code{05\_protein\_properties.py} & Figure 5 descriptors \\
\code{06\_motif\_analysis.py} & Figure 7 motifs \\
\code{07\_novelty\_audit.py} & the 62.3\% verbatim-copy rate \\
\code{08\_extract\_table\_s1.py} & the paper's own sequences, MW-verified \\
\code{09\_verify\_published\_sequences.py} & the exact-match result \\
\code{09b\_verify\_against\_tables\_s2\_s3.py} & descriptor and novelty cross-checks \\
\code{10\_ablation\_bestofn.py} & best-of-$N$ curves \\
\code{10b\_scoring\_stochasticity.py} & the rejected explanation for the gap \\
\code{11\_ablation\_sequence\_content.py} & the sequence ablations \\
\code{12\_baseline\_composition.py} & the composition baselines \\
\code{13\_extension\_nonmasp.py} & out-of-distribution families \\
\code{14\_composition\_association.py} & composition association + null \\
\code{99\_make\_figures.py} & all six figures, from saved results \\
\bottomrule
\end{longtable}

Every script writes a JSON manifest under \code{results/} recording the git
commit, working-tree cleanliness, model revision, seed, platform, GPU and the
version of every library used.

\section{Glossary}

\begin{description}[leftmargin=0cm,style=nextline,itemsep=3pt]
\item[Amino acid] One link in a protein chain; twenty kinds, written as letters.
\item[Attention] The mechanism by which each position in a transformer draws on
  every other, with learned, content-dependent weights.
\item[Autoregressive] Generating one token at a time, each conditioned on all
  previous ones.
\item[Best-of-$N$] Generating $N$ candidates and keeping the highest-scoring.
  The expected result rises with $N$ regardless of model quality.
\item[BPE] Byte-pair encoding; builds a vocabulary by repeatedly merging the
  most frequent adjacent pair.
\item[Decoding] Turning next-token probabilities into an actual choice: greedy,
  or sampling with temperature / top-$k$ / top-$p$.
\item[Dragline] The silk a spider hangs from; made of MaSp; the silk that gets
  mechanically tested.
\item[Fine-tuning] Further training of a pretrained model on a specific task.
\item[Held-out] Data withheld from training, used to measure generalisation.
  There is none here.
\item[In-sample] Measured on training data. An upper bound, not an estimate.
\item[Instability index] A ProtParam statistic computed from dipeptide
  frequencies; estimates whether a protein is stable.
\item[Isoelectric point] The pH at which a protein carries no net charge.
\item[MaSp] Major ampullate spidroin; the dragline protein; the paper's subject.
\item[Motif] A short recurring sequence pattern, e.g.\ poly-A, GGX, GPGXX.
\item[Normalisation] Mapping raw measurements to $[0,1]$ by min--max scaling.
\item[Parameter] One learned number in a model. 253.6 million here.
\item[pLDDT] AlphaFold2's per-residue confidence score. 40--60 is low.
\item[Poly-A] A run of alanines; forms $\beta$-sheet nanocrystals; supplies
  strength and stiffness.
\item[ProtParam] A standard tool computing descriptors from a sequence.
\item[\Rsq] $1 - \mathrm{SS}_{\text{res}}/\mathrm{SS}_{\text{tot}}$. Negative
  values mean worse than guessing the mean.
\item[RoPE] Rotary positional embedding; encodes position by rotation, so
  relative offsets are what the attention sees.
\item[Spidroin] A spider silk protein.
\item[Temperature] A decoding parameter rescaling the distribution; lower is
  more deterministic.
\item[Token] The unit a model consumes; here, a chunk of amino acids.
\item[Toughness] Energy absorbed per unit volume before breaking; the area
  under the stress--strain curve.
\item[Within-vector \Rsq] The paper's definition: \Rsq{} across the eight
  property entries of one sequence.
\end{description}

\section{Reproducing the reproduction}

\begin{lstlisting}[style=plain]
python -m venv .venv && .venv/Scripts/activate
pip install torch --index-url https://download.pytorch.org/whl/cu126
pip install -e ".[dev]"

git clone https://github.com/lamm-mit/SilkomeGPT _ref/SilkomeGPT-main
# plus the two Silkome files; see data/raw/README.md

python scripts/00_check_env.py
python scripts/03_build_dataset.py --fasta <fasta> --mech <csv>
python scripts/02_reproduce_table1.py --n 2048     # ~4.5 h on a 6 GB GPU
python scripts/99_make_figures.py
\end{lstlisting}

Two things deliberately do not ship: the Silkome primary data and the authors'
checkpoint, neither of which is mine to redistribute. Both are freely
available, and \code{data/raw/README.md} says exactly what to fetch.

\section*{Sources}
\addcontentsline{toc}{section}{Sources}

\begin{enumerate}[leftmargin=*,itemsep=2pt]
\item W.~Lu, D.~L.~Kaplan, M.~J.~Buehler. Generative Modeling, Design, and
  Analysis of Spider Silk Protein Sequences for Enhanced Mechanical Properties.
  \emph{Adv.\ Funct.\ Mater.} \textbf{34}, 2311324 (2024).
  \href{https://doi.org/10.1002/adfm.202311324}{doi:10.1002/adfm.202311324}.
  Preprint with the supplementary tables:
  \href{https://arxiv.org/abs/2309.10170}{arXiv:2309.10170}.
\item K.~Arakawa \emph{et al.} 1000 spider silkomes: linking sequences to silk
  physical properties. \emph{Sci.\ Adv.} \textbf{8}, eabo6043 (2022).
\item A.~Pandey, W.~Chen, S.~Keten. Primary sequence-based data-constrained
  machine learning framework to predict mechanical properties of the spider
  silk. \emph{Commun.\ Mater.} (2024). \emph{(Source of the composition
  analysis in §\ref{sec:nullretraction}, not of anything in Lu et al.)}
\item Released model and notebook:
  \href{https://github.com/lamm-mit/SilkomeGPT}{github.com/lamm-mit/SilkomeGPT},
  \href{https://huggingface.co/lamm-mit/SilkomeGPT}{huggingface.co/lamm-mit/SilkomeGPT}.
\end{enumerate}
